%============================================================
%  LEX v3.0 — Template Artikel Jurnal Hukum
%  Standard: Jurnal Konstitusi / Sinta 1 / Scopus Q1
%  Gunakan /draft-subbab di Copilot Chat untuk mengisi bagian ini
%============================================================

\documentclass[12pt,a4paper]{article}

%--- Paket utama ---
\usepackage[bahasa]{babel}
\usepackage[utf8]{inputenc}
\usepackage[T1]{fontenc}
\usepackage{times}                % Font Times New Roman
\usepackage[margin=2.5cm]{geometry}
\usepackage{setspace}
\doublespacing

%--- Sitasi & Bibliografi ---
\usepackage{hyperref}
\hypersetup{colorlinks=true, linkcolor=black, citecolor=black, urlcolor=blue}

%--- Visualisasi Akademik (LEX v3.0) ---
\usepackage{tikz}
\usetikzlibrary{shapes.geometric, arrows.meta, positioning, fit, backgrounds}
\usepackage{pgfplots}
\pgfplotsset{compat=1.18}
\usepackage{pgf-pie}
\usepackage{booktabs}
\usepackage{multirow}
\usepackage{array}

%--- Format ---
\usepackage{titlesec}
\usepackage{abstract}
\usepackage{enumitem}
\usepackage{footnote}
\makesavenoteenv{tabular}

%============================================================
%  METADATA ARTIKEL
%  Isi bagian ini terlebih dahulu sebelum menulis
%============================================================
\title{%
  \textbf{[JUDUL ARTIKEL]}\\[0.5em]
  \large [Subjudul jika ada]
}

\author{%
  [Nama Penulis Pertama]\textsuperscript{1} \and
  [Nama Penulis Kedua]\textsuperscript{2}\\[0.5em]
  \small\textsuperscript{1}[Institusi, Email]\\
  \small\textsuperscript{2}[Institusi, Email]
}

\date{Diterima: [tanggal] \quad Disetujui: [tanggal]}

%============================================================
\begin{document}
\maketitle

%------------------------------------------------------------
\begin{abstract}
\noindent
% Perintah Copilot: /draft-subbab Abstrak
% Target: 150-250 kata, bahasa Indonesia dan Inggris
[Tulis abstrak di sini. Copilot akan menghasilkan teks setelah Anda menggunakan perintah /draft-subbab Abstrak]

\medskip
\noindent\textbf{Kata kunci:} [kata kunci 1]; [kata kunci 2]; [kata kunci 3]; [kata kunci 4]; [kata kunci 5]

\medskip
\noindent\textbf{Abstract:} [English abstract]

\noindent\textbf{Keywords:} [keyword 1]; [keyword 2]; [keyword 3]; [keyword 4]; [keyword 5]
\end{abstract}

%------------------------------------------------------------
\section{Pendahuluan}
% Perintah Copilot: /draft-subbab Pendahuluan
% Target: 1000-1500 kata
% Struktur: latar belakang → rumusan masalah → tujuan → metode → sistematika
% ATD: Norma Positif + Asas + Teori + Doktrin

[Copilot akan mengisi bagian ini dengan perintah /draft-subbab Pendahuluan]

%------------------------------------------------------------
\section{Metode Penelitian}
% Perintah Copilot: /draft-subbab Metode
% Jenis penelitian, pendekatan, bahan hukum, teknik analisis

[Copilot akan mengisi bagian ini dengan perintah /draft-subbab Metode]

%------------------------------------------------------------
\section{Pembahasan}

\subsection{[Judul Subbab A]}
% Perintah Copilot: /draft-subbab Pembahasan A
% ATD WAJIB: Asas + Teori + Doktrin terintegrasi di setiap paragraf

[Copilot akan mengisi bagian ini]

\subsection{[Judul Subbab B]}
% Perintah Copilot: /draft-subbab Pembahasan B
% Sisipkan visualisasi dari templates/atd-diagram.tex atau templates/visualisasi.tex

[Copilot akan mengisi bagian ini]

\subsection{[Judul Subbab C]}
% Perintah Copilot: /draft-subbab Pembahasan C

[Copilot akan mengisi bagian ini]

%------------------------------------------------------------
\section{Penutup}

\subsection{Kesimpulan}
% Perintah Copilot: /draft-subbab Kesimpulan
% Jawaban langsung atas rumusan masalah, tanpa pengulangan pembahasan

[Copilot akan mengisi bagian ini]

\subsection{Saran}
% Perintah Copilot: /draft-subbab Saran
% Rekomendasi kebijakan/legislasi konkret

[Copilot akan mengisi bagian ini]

%------------------------------------------------------------
\bibliographystyle{apalike}
\bibliography{../referensi/pustaka}

\end{document}
